% !TeX root = ../车云协同液罐车防侧翻行驶规划与控制研究.tex
\chapter{总结与展望}

\section{全文总结}

液罐车侧翻事故具有诱因耦合强、危险演化快、后果严重等特点。面向高等级公路典型运行场景，传统单车智能方法在感知范围、先验信息获取能力以及对长时域交互风险的提前规避能力方面存在天然局限；同时，液体晃动与车辆横纵向动力学之间的强耦合作用，又使液罐车在紧急变道、避障和高速转向等工况下表现出显著区别于普通车辆的稳定性约束特征。围绕上述问题，本文以车路云一体化系统为背景，系统开展了车云协同液罐车防侧翻行驶架构建模、预测性决策与防侧翻规划、车端轨迹跟踪抑晃防侧翻控制以及仿真与模型试验验证研究，形成了“体系架构—云端决策规划—车端控制保障—多层验证”的完整技术链路。通过上述研究，本文得到如下主要结论。

\begin{enumerate}[itemsep=2pt, parsep=0pt]
    \item 针对液罐车防侧翻问题中长期缺乏体系层设计依据、难以统一组织云端预测规避与车端实时控制的问题，本文提出了面向车云协同液罐车防侧翻场景的分场景基于模型与仿真的体系工程建模方法，构建了CT2ROP场景任务体系架构，形成了“战略—运行—服务—资源”四阶段递进的架构建模与闭环验证流程。该方法从IVCPS开放演进、多主体协作和能力跨层传递的特征出发，将液罐车防侧翻场景中的能力目标、运行主体、服务链路和资源支撑统一纳入同一建模框架，并通过能力测度与蒙特卡洛仿真对关键服务指标进行了验证，从而为后续算法研究提供了清晰的系统边界、时间尺度约束与接口组织依据。
    
    \item 针对复杂高速交通场景中周车行为具有显著交互性和非确定性、传统短视行为决策难以提前规避液罐车潜在侧翻诱因的问题，本文构建了图扩散交互式预测模型与强化学习长时域决策模型，实现了云端面向长时域风险演化的预测性决策。该方法以像素级道路有向交通图表征场景结构，通过自回归图扩散流匹配学习周车未来演化分布，获得对周车多模态交互行为的长时域预测；进一步结合车辆参数与交互意图识别结果训练强化学习决策器，输出速度与换道行为决策，并通过防侧翻换道轨迹生成方法构造满足平滑性与侧翻约束要求的参考轨迹。该研究增强了系统对诱发性风险的前瞻识别与提前规避能力，为液罐车在复杂交互环境中的安全决策提供了云端支撑。
    
    \item 针对云端预测性决策结果在局部突发场景下仍需要车端快速重构安全轨迹的问题，本文提出了车端防侧翻避障应急轨迹规划方法，建立了面向半挂液罐车的运动学、运动动力学耦合及简化侧向动力学模型，构造了约束迭代LQR求解框架。该方法将障碍避让约束、道路边界约束、车辆动力学约束和防侧翻约束统一纳入局部规划问题，并结合关键周车筛选与双场景应急避障策略，实现了在紧急工况下的快速局部轨迹重构。仿真结果表明，所提方法能够在车端形成面向局部风险的快速响应能力，与云端长时域决策共同构成“前瞻规避—局部重规划”的分层规划链。
    
    \item 针对液罐车液体晃动与车辆横纵向运动耦合显著、现有快速模型难以同时兼顾实时性与对关键动力学特征的刻画能力的问题，本文提出了横纵向耦合液体晃动滑动伸缩天平模型及残差强化学习重线性化方法。所建滑动伸缩天平模型在快速计算框架下表征了液体质量转移、横纵向力与附加力矩等关键作用机制；进一步通过线性双摆模型与残差强化学习重线性化模型相结合，实现了对高保真非线性液体晃动动力学的在线逼近。该研究为后续控制器设计提供了兼具计算效率和精度支撑的预测模型，使液罐车防侧翻控制能够在保留关键车液耦合特征的基础上满足实时应用需求。
    
    \item 针对液罐车在极限转向与避障工况下需要同时满足轨迹跟踪、抑晃与防侧翻多目标要求的问题，本文设计了多约束模型预测控制策略，建立了横纵向耦合半挂液罐车动力学模型，并将等效载荷转移率、执行器约束和舒适性约束统一纳入优化控制问题。该方法通过对侧翻风险指标和车辆动态响应的协同约束，实现了液罐车在复杂参考轨迹下的稳定跟踪与姿态调节；同时结合差动制动和工程补偿机制，提高了控制器在高风险工况下的稳定保持能力。该研究形成了面向液罐车防侧翻的车端最终安全保障方法，使系统在进入危险状态邻域后仍能够维持可控。
    
    \item 搭建了车液耦合联合仿真平台与缩比例模型车试验平台，对所提车云协同液罐车防侧翻方法进行了系统验证。仿真层面围绕左转弯、高速单移线、短距双移线及麋鹿工况等典型高风险场景开展了分析，验证了云端决策规划方法和车端控制方法在风险规避、姿态抑制与轨迹跟踪方面的有效性；试验层面基于缩比相似性准则构建了模型车硬件与自动驾驶软件系统，对车端控制策略进行了实证检验。结果表明，本文所提出的方法能够在多类复杂工况下有效降低液罐车侧翻风险，提升轨迹跟踪性能与姿态稳定性，并具备较好的工程实现可行性。
\end{enumerate}

\section{主要创新点}

本文的主要创新点可归纳为以下三个方面。

\begin{enumerate}[itemsep=2pt, parsep=0pt]
    \item 创建了车云协同液罐车防侧翻场景任务体系架构。针对液罐车防侧翻研究中长期存在的“架构层缺位、方法层孤立”的问题，本文进行了战略、运行、服务、资源四阶段一体化建模，并通过仿真验证完成架构闭环，为后续算法模块的问题边界、接口衔接、时间尺度等参数划分及工程部署提供了可追溯的分析依据。
    
    \item 提出了面向交通交互场景的图扩散云端预测性决策与规划方法。针对液罐车在高速场景下面临的周车交互不确定性强、长时域风险演化建模难、现有表征方法信息冗余大等问题，本文提出交通图表征方法、自回归图扩散交互式预测模型及强化学习长时域决策模型，建立了从周车未来行为预测到液罐车横纵向耦合决策再到防侧翻平滑参考轨迹生成的完整方法链，实现了面向诱发性风险的前瞻规避。
    
    \item 提出了面向液罐车横纵向耦合晃动特性的抑晃防侧翻轨迹跟踪控制方法。针对液罐车车液耦合动力学计算复杂及现有研究仅考虑横向动力学的问题，本文构建了高精度反映罐液耦合动态的滑动伸缩天平快速计算模型及其残差强化学习重线性化方法，并基于此设计了多约束模型预测控制器以抑制状态风险进一步演化，形成了与长时域前瞻规划协同的短时域高频安全保障。
\end{enumerate}

\section{研究展望}

尽管本文围绕车云协同液罐车防侧翻行驶规划与控制开展了较为系统的研究，并取得了阶段性成果，但仍有若干问题值得进一步深入研究。未来可从以下几个方面继续拓展。

\begin{enumerate}[itemsep=2pt, parsep=0pt]
    \item 面向更复杂交通场景的模型扩展与泛化研究。本文当前研究对象主要聚焦于高等级公路典型工况，后续可进一步拓展至匝道汇入、连续弯道路段、施工区域以及城市快速路等更复杂场景，并引入更加丰富的驾驶风格、交通参与体行为模式和环境扰动因素，以增强云端交互预测、行为决策与车端局部规划方法的泛化能力。
    
    \item 面向通信不确定性与系统安全性的协同增强研究。车云协同系统的实际应用仍受到通信时延、丢包、带宽波动及信息安全等因素影响。后续可进一步研究通信不确定条件下的规划与控制平滑切换机制、任务降级机制以及轻量化安全认证机制，增强车云协同液罐车系统在复杂网络条件下的鲁棒性、可信性与可部署性。
    
    \item 面向真实工程应用的软硬件集成与实车验证研究。本文已通过联合仿真与缩比例模型试验验证了方法有效性，后续可进一步结合真实液罐车平台开展道路试验，积累多源运行数据并进行闭环迭代优化；同时可探索与数字孪生、高保真世界模型以及更强表达能力学习框架的融合，进一步提升车云协同液罐车防侧翻系统的实用化与智能化水平。
\end{enumerate}
